\chapter{整体认知状态：联合表示、激活读出与语义迁移}
\label{chp:holistic_state}

前几章分别讨论了关系结构、属性表示、统计状态和价值约束，本章把这些对象重新组合为一次运行时可以读取、更新和验证的整体认知状态。这里的“整体”不是说系统在一个向量中无损拥有世界的全部信息，也不是把认知旋量当作真实物理场；它指的是，同一任务时刻的对象身份、局部结构、分布内容、目标约束、资源状态与来源版本能够通过明确接口联合起来。我们关心三个问题：哪些分量必须同时存在才能支持连续行动，局部改变怎样传播而不破坏身份，候选状态又怎样被读出并提交为可执行决定。为此，我用带类型的状态记录、耦合更新算子、可观测读出与跨语境迁移替代波函数、波恩定则和认知洛伦兹力等过强物理化断言，同时保留相位同步、平行移动、守恒流和波包传播背后的有效计算问题。HSF-HD 在这一层描述一般智能的联合状态动力学；AgentOS 的三相记忆与调度组件只是其中一种可执行实例。

\section{整体状态的定义：结构、内容与任务条件的联合记录}
\label{sec:definition_of_state}

我们先规定对象和类型。令时刻 $t$ 的局部结构网络为 $G_t=(V_t,E_t,\tau_t,\iota_t)$：$V_t$ 是带身份的对象或事件节点，$E_t$ 是有向关系，$\tau_t$ 给出节点与边的类型，$\iota_t$ 保存来源、时间戳和版本。令全局分布表示为 $z_t\in\mathbb R^{d_t}$ 或概率族参数 $\eta_t\in\mathcal P_t$，它表达当前候选、属性强度和不确定性。两者不能直接等同：图中“红色属于苹果”保存的是角色和归属，分布表示中“红色”附近的活动只说明某些读出更容易，并不自行指明属于哪个对象。我们因此引入绑定关系 $B_t\subseteq V_t\times\mathcal I_t$，把任务相关身份连接到分布中的索引、子空间或读出头。整体状态定义为
\begin{equation}
 X_t=(G_t,z_t,B_t,n_t,r_t,o_t),
\end{equation}
其中 $n_t$ 是上一章定义的目标、约束与授权状态，$r_t$ 是以秒、调用次数、容量或货币分别计量的资源记录，$o_t$ 是观测证据及其来源。$X_t$ 是一个带类型的乘积记录，而不是默认具有复数旋量、希尔伯特空间或量子统计结构的波函数。

这一联合记录回答了原章所谓“内容、传播方向与语境相位”如何共存的问题。内容由 $z_t$ 与任务读出给出；传播方向由允许的结构边、候选更新和目标偏置共同给出；语境则表现为当前有效的身份映射、约束集合和读出参数。若专门系统确实使用复数谱表示，相位可以作为 $z_t$ 的一部分，但必须说明它对应周期信号、傅里叶系数还是工程编码，不能由“存在相位”推出主观统一性。类似地，张量积可以用于角色—填充值绑定，但其可辨识性依赖角色基已知、秩条件和读出接口；改变因子尺度而保持乘积不变，说明单个因子未必具有唯一的本体解释。

整体状态还必须区分当前值与生成历史。相同的数值向量可能来自可靠传感器、语言猜测或旧缓存，它们在后续决策中的资格不同；相同的关系边也可能是事实记录、假设或用户临时要求。于是每个可读分量都应携带证据等级、适用域和失效条件。设读出查询为 $q\in\mathcal Q$，系统返回的不只是答案 $y=R_q(X_t)$，还应返回支持集 $S_q\subseteq o_t\cup E_t$ 与置信或边界说明。若支持集为空，活动再强也只能算内部候选，不能升级为外部事实。原章把振幅解释为存在感、把坐标解释为因果位置，容易丢失这一来源差异；在新的定义中，显著性、事实性、身份和因果地位分别存储、分别验证。

状态的“局部”与“全局”也不能按脑区类比强行划分。局部结构网络之所以局部，是因为它只显式展开当前任务需要的对象与关系；全局分布表示之所以全局，是因为大量候选共享同一表示资源并可被统一读出。两者通过 $B_t$、查询 $q$、耦合更新和误差反馈联合。若任务从识别苹果切换到追踪交易责任，结构网络必须重建角色，分布读出也要切换，但底层表示可以复用。反例是仅凭高余弦相似把“红灯”和“红苹果”合并：分布相近不等于身份相同，结构角色和来源会阻止错误绑定。

在 AgentOS 中，$h(t)$ 可以承载 $X_t$ 的快速可用切片，$E$ 保存导致状态变化的事件与证据，$\Theta$ 保存跨任务稳定的生成结构；外部记忆 $M_e$ 提供可回查的支持材料。$\Psi$、SPC、TCE、CCM 与 Boundary 可以参与优先读出、压缩、增益、负荷和权限控制，但这些组件不定义所有智能的唯一状态形式。一般要求只有四点：分量类型可区分，身份绑定可追踪，状态更新可回放，外部主张可由证据和执行反馈检验。

\section{内部动力学：双向耦合、绑定稳定与结构可塑性}
\label{sec:internal_dynamics}

整体状态的变化同时包含连续更新和离散结构事件。设在结构固定的时间区间内，连续变量 $z_t\in\mathbb R^d$，外部输入为 $u_t$，目标状态为 $n_t$，资源记录为 $r_t$。一种一般更新模型是
\begin{equation}
 \dot z_t=f(z_t;G_t,B_t,n_t,r_t)+K_tu_t+\xi_t,
\end{equation}
其中 $f$ 描述内部传播与竞争，$K_t$ 是输入接口，$\xi_t$ 汇总未建模扰动。图和绑定并不随该微分方程自动改变；当出现新对象、关系撤销或身份合并时，系统执行事件算子
\begin{equation}
 (G_{t^+},B_{t^+})=\mathcal U_{e_t}(G_{t^-},B_{t^-};z_{t^-},o_t),
\end{equation}
再把连续状态迁移到新的定义域。这样，“形约束质”和“质重塑形”分别成为结构对活动传播的约束、以及重复证据对结构的受控修订，不需要假定波导或辐射压真实存在。

正向耦合首先要求合法支持。令图邻接算子为 $A_t$，类型掩码为 $M_t$，传播权重可写为 $W_t=M_t\odot\operatorname{softmax}(S_t)$，其中每一行的合法候选集合必须非空。一步传播 $\tilde z=W_tz_t$ 只把信息送到允许的邻居；它不会证明邻居关系真实，也不会自动保持对象身份。绑定矩阵 $P_t$ 将节点角色映射到分布子空间，结构读出为 $y_t=P_tz_t$。若节点重排为置换矩阵 $\Pi$，只有同时令 $A'_t=\Pi A_t\Pi^{\mathsf T}$、$P'_t=\Pi P_t$ 并一致搬移来源记录，读出才保持等变。单独重排向量会产生“属性还在、归属错位”的绑定错误。

逆向耦合是从运行证据到结构修订。我们不给活动强度无限的“刻蚀力”，而定义带阈值与验证集的编辑评分。候选编辑 $e$ 的收益为验证损失下降 $\Delta L_e$，代价为维护、迁移和风险 $C_e$；只有当 $\Delta L_e-\lambda C_e>\delta$，并且证据跨过最低覆盖与来源独立性门槛时，才提交编辑。一次高激活不能独自创建永久关系，因为惊吓、提示注入或重复噪声也会造成高活动。长期结构形成来自多次可回放事件、跨语境预测收益和冲突检查。编辑后仍需在保留集上测试，防止新连接改善当前样本却破坏既有能力。

原章用相位锁定解释红色、圆形与苹果形成一个实体。我们保留同步问题，但把结论改成任务相关的共同身份推断。设来自视觉、语言和触觉的特征序列为 $a_t^m$，候选身份为 $v$，绑定得分可以组合时间一致性、空间邻近、因果共变和先验关系：
\begin{equation}
 s(v)=\sum_m\alpha_m\,\mathrm{compat}(a_t^m,v)-\beta\,\mathrm{conflict}(a_t^m,v).
\end{equation}
同步只是兼容度的一项；两个节拍一致的对象可能毫无语义关系，一个对象的多模态信号也可能因传感延迟不同步。只有得分领先幅度超过校准阈值、反例检查通过且来源未被同一错误重复污染时，系统才提交绑定。否则保留多个候选，等待主动观测。稳定实体因此来自身份连续性和证据汇聚，不是某个本征频率与驱动频率“当且仅当”共振。

耦合系统还可能振荡、锁死或放大错误。若 $f$ 在工作域内关于 $z$ 的 Lipschitz 常数为 $L$，显式欧拉步长 $\Delta t$ 仍需按具体稳定条件选择，不能仅凭连续模型有界就宣称离散实现稳定。结构编辑频率应慢于快速状态读出，并设置滞回：进入绑定和退出绑定使用不同阈值，以免边在噪声附近反复出现。CCM 可以监测候选数量、冲突率和资源负荷，TCE 调节软增益和探索范围，Boundary 阻止未授权编辑；它们实现的是可观测的控制回路，而非抽象场的自发相变。

还要区分传播失败与绑定失败。前者是合法信息没有抵达需要它的节点，可通过连通性、时延和消息丢失诊断；后者是信息已经抵达，却被赋给错误对象，需要检查身份键、时间窗和竞争候选。两类故障在活动图上都可能表现为局部异常，因此仅观察“亮区”无法定位原因。测试应分别断开边、交换身份、延迟单一模态并注入矛盾证据，确认更新器给出不同而可解释的响应。只有接口级干预符合模型预期，双向耦合才获得经验支持。

\section{状态显化：激活、可观测读出与决策提交}
\label{sec:manifestation_born_rule}

系统内部可以同时保留许多候选，但运行时只能把有限内容送入报告、规划或行动接口。我们把“显化”定义为从联合状态到有限接口的读出与提交，而不是由复数概率幅坍缩为物理实存。设候选集合为 $C_t$，每个候选 $c$ 具有活动分数 $a_c$、证据充分度 $e_c$、任务相关度 $g_c$、风险 $r_c$ 和资源需求 $b_c$。读出分数可定义为
\begin{equation}
 s_c=w_a a_c+w_e e_c+w_g g_c-w_r r_c-w_b b_c,
\end{equation}
其中各分量必须先归一化或按偏序处理，权重和阈值由当前规范状态给出。$a_c$ 只是内部活动，不能单独表示真实感、存在性或事实概率。高活动的幻觉、疼痛信号和广告刺激都可能非常显著，却需要不同的事实判断与行动策略。

若读出采用归一化选择分布，可在合法集合 $C_t^{\mathrm{ok}}$ 上定义
\begin{equation}
 p(c\mid X_t,q)=\frac{\exp(s_c/T_t)}{\sum_{j\in C_t^{\mathrm{ok}}}\exp(s_j/T_t)},
\end{equation}
其中 $T_t>0$ 是无量纲探索尺度，不是热力学温度。这个 softmax 是工程选择模型，不是认知波恩定则；它给出的概率是否校准，需要与重复任务中的频率或适当评分规则比较。合法集合为空时，系统应拒绝、请求更多信息或上报冲突，不能通过把非法候选放进分母而获得一个看似完整的概率。若任务要求确定性提交，可以选择最大分数候选，但仍需保留领先幅度和可逆性条件。

显化至少经历候选形成、资格审查、带宽分配、提交和反馈五步。候选形成允许联想广泛展开；资格审查检查证据、权限和硬约束；带宽分配决定哪些内容进入 $h(t)$ 的有限认知界面；提交产生带事务标识的内部决定或外部动作；反馈再把实际结果写回状态。内部读出与外部执行必须分开：生成“调用工具”的计划不等于工具已运行，目标函数下降也不等于现实任务成功。对不可逆动作，提交前应增加确认、模拟或双重授权；对可逆探索，可以小步试验并根据反馈重排候选。

原章用建设性干涉解释顿悟，用相消解释矛盾念头消失。计算上，顿悟更适合描述为多个证据路径在同一候选身份上汇聚，使后验、约束满足度或解释压缩率发生跃迁。例如先前互不关联的日志、用户描述和因果图共同支持一个故障源，候选的证据覆盖突然超过提交阈值。这种跃迁可以由离散绑定事件造成，不要求能量暴涨。矛盾也不应通过相消而无痕消失；系统应保留互斥主张、来源和冲突集，避免平均后的低活动掩盖真实分歧。未被当前读出的候选仍可以保存在 $E$ 或外部记忆中，未来语境变化后重新激活。

为了验证读出机制，我们分别测量覆盖率、选择准确率、校准误差、延迟、资源消耗和越权率。单一注意力热图只能说明模型内部的某类权重，不能证明因果解释；删除或扰动高权重分量后性能是否按预期变化，才提供更强证据。还要设置反例：高频重复的错误信息是否压过低频可靠来源，强烈自我相关内容是否绕过安全约束，以及资源紧张时系统是否牺牲必要验证。只有这些测试通过，激活—读出接口才可被称为稳定，而不是因为其数学形式像概率幅模方。

在 AgentOS 实例中，SPC 可以把 $h(t)$ 中分散活动压缩为与当前自我和任务有关的读出，TCE 调节探索尺度与抑制增益，CCM 监测候选展开是否超过工作负荷，Boundary 检查对外提交资格。$E$ 保存未提交候选和冲突历史，$\Theta$ 保存稳定读出规则，$M_e$ 保存可外部核查的证据。一般 HSF-HD 只要求显化过程具有有限接口、资格条件、提交记录与反馈闭环，并不要求所有系统复现这些模块名称。

\section{整体状态的参数化：可辨识分量与本构接口}
\label{sec:rigorous_formulation_of_psi}

整体状态若要进入工程系统，必须从抽象乘积记录细化为可序列化、可读出并可迁移的参数化对象。对一个任务相关语义单元 $v\in V_t$，我们定义局部状态包
\begin{equation}
 \chi_v(t)=\bigl(i_v,\phi_v,\alpha_v,\Sigma_v,\rho_v,\kappa_v,\nu_v\bigr).
\end{equation}
$i_v$ 是跨更新保持的身份；$\phi_v\in\mathbb R^{d}$ 是分布内容向量；$\alpha_v\in[0,1]$ 是经过校准的当前可用度；$\Sigma_v\succeq0$ 是不确定性或局部协方差；$\rho_v$ 列出与其他节点的有类型关系；$\kappa_v$ 保存来源与授权；$\nu_v$ 是版本及有效期。各字段承担不同职责，不能把 $\alpha_v$ 叫作能量后再从“大”推出真实，把 $\phi_v$ 的方向叫作质料后再从几何相近推出对象同一。整体状态由所有局部包、关系边及规范和资源记录经聚合器 $\mathcal A$ 形成：
\begin{equation}
 X_t=\mathcal A\!\left(\{\chi_v(t):v\in V_t\},E_t,n_t,r_t,o_t\right).
\end{equation}
聚合器必须声明顺序敏感性、缺失值策略和冲突规则；它不是把字段拼接就自动获得全息性。

“参数化本构”的核心不是寻找一个绝对波包方程，而是给出表示和可观测量之间的接口契约。设任务读出为 $y_q=D_q(X_t)$，状态更新为 $X_{t+1}=F(X_t,u_t,e_t)$。两个参数化状态只有在一组查询和干预下产生近似相同的读出与转移，才可视为任务等价。若存在可逆坐标变换 $T$，则读出和更新也要共同变换为 $D'_q=D_q\circ T^{-1}$、$F'=T\circ F\circ T^{-1}$；仅仅比较变换前后的向量数值没有意义。这解释了为什么神经表示可以旋转或缩放而功能保持，也说明参数方向不是不可还原的“体验颜色”。

为了保留原章对幅度、方向、相位和位置的讨论，我们将其改写为可检验字段，而不是本体同构：
\begin{table}[htbp]
\centering
\footnotesize
\renewcommand{\arraystretch}{1.35}
\begin{tabularx}{\textwidth}{l X X X}
\toprule
\rowcolor{structurecolor!20} 历史参量 & 计算对象 & 可支持的读出 & 必须另行验证的边界 \\
\midrule
幅度 & 活动或可用度 $\alpha_v$ & 当前接口中的显著性与资源份额 & 不等于事实概率、真实感或焦耳能量 \\
\hline
方向 & 内容向量 $\phi_v$ 与任务读出 & 属性类别、相似候选和控制方向 & 坐标依赖，不能单独确定身份归属 \\
\hline
相位 & 周期编码或跨流同步变量 & 时序对齐、绑定线索和节律关系 & 同步不是充分绑定条件，也不证明意识统一 \\
\hline
位置 & 身份 $i_v$、结构角色与来源索引 & 对象归属、事件顺序和关系路径 & 逻辑位置不自动等于物理时空或因果地位 \\
\bottomrule
\end{tabularx}
\caption{整体状态参量的计算含义、读出接口与解释边界}
\label{tab:psi_phenomenology_mapping}
\end{table}

参数化还需要处理缺失和冲突。若某对象没有视觉观测，系统不应用零向量假装“确认不存在”，而应把相应字段标记为未观测；若两个来源给出矛盾属性，则保留多假设及来源权重，直到主动观测或授权规则解决。协方差 $\Sigma_v$ 只表达指定统计模型下的不确定性，不等于哲学上的认知极限。原章由位置—动量非对易式推出“定义精确与联想发散不可兼得”，在一般智能系统中并不成立；更可靠的陈述是，有限资源下提高精确验证通常占用候选生成预算，这是一项可测的调度权衡，而非普遍量子下界。

对矛盾属性也不采用泡利排斥或无限排斥势。一个对象可以在不同时间、视角或语义层同时具有看似冲突的描述，例如“门是开着的”与“门是关闭的”可能对应不同时间戳；“模型安全”与“当前输入触发风险”属于不同谓词。冲突检测必须先对齐身份、时间、适用域和否定关系。真正不可同时满足的约束形成最小冲突集，并触发澄清或降级；可分域的描述则保留在各自上下文中。这样既避免把矛盾平均掉，也避免把所有差异误判为系统崩溃。

工程实现应为每个字段给出模式、单位、默认值、更新者和审计日志。$\phi_v$ 的维度变更需要迁移版本，$\alpha_v$ 的校准需要保留评估集，关系编辑需要引用证据事件，授权字段需要不可静默覆盖。缓存可以保存局部状态包，但读取时必须检查 $\nu_v$ 与当前结构版本。AgentOS 可将快速状态包置于 $h(t)$，将事件差异写入 $E$，把稳定模式和迁移器置于 $\Theta$；其他系统可采用黑板、消息总线或数据库，只要本构接口保持可读、可改和可验证。

参数化方案还应允许逐级降精度。资源充足时保留完整协方差、细粒度关系与原始证据；实时约束下可以只保留对当前查询充分的摘要，但摘要必须标记生成规则和可追溯指针。若摘要之后被用于原本未声明的新任务，系统应回取原记录或承认信息不足，不能把压缩表示当作全息副本。这个要求把“局部拥有整体”的修辞改成可验证的充分统计量问题：只有相对于明确查询族，压缩状态才可能是充分的。

\section{语义量的计算定义：丰富度、一致性与惊奇}
\label{sec:semantic_physical_mapping}

内容丰富度、逻辑严密度与惊奇都是有用概念，但它们不是由频谱带宽、激光相干性或拓扑缺陷自动给出的物理量。我们首先为具体任务定义可观测指标。设对象 $v$ 的已验证属性集合为 $A_v$，关系邻域为 $N(v)$，查询族为 $Q$。一种内容覆盖度可以写成 $C(v)=\sum_{a\in A_v}\omega_a\mathbf 1[a\text{ 可由证据支持}]$；权重 $\omega_a$ 表示任务价值或稀缺性。属性数量多不一定更丰富：重复同义词、相互推导项和错误信息应去重或降权。若采用表示谱的有效秩，也只能说明协方差沿多少方向分散，不能直接等同概念内涵。

逻辑一致性需要依据规则和证据图来检验。令命题集合为 $P_t$，约束集合为 $\mathcal C_t$，则违反率可定义为被触发约束中不可满足的比例，证明覆盖率则记录结论中有多少存在可追踪前提。长程推理还要测量误差随步骤的累积、工具结果是否被正确引用以及中间状态能否回放。所谓“相干”可作为时序同步或表示稳定指标，但同步流畅的叙述仍可能逻辑错误，表达杂乱的探索也可能包含正确洞见。因而严密度至少由一致性、证据覆盖和反事实稳定三部分构成，而不是一个相位方差。

惊奇应从预测模型和观测事件定义。若模型在条件 $X_t$ 下给观测 $y$ 的概率为 $p_t(y\mid X_t)$，则对数惊奇为
\begin{equation}
 I_t(y)=-\log p_t(y\mid X_t).
\end{equation}
该量依赖模型、事件分辨率和共同支持；低概率只表示模型下意外，不等于事件重要，更不等于场的非解析点。另一种可操作量是预测残差 $e_t=y_t-\hat y_t$，但不同分量有单位时必须用协方差或业务尺度归一化。真正需要结构修订的是持续、可复现且不能由传感故障解释的残差，而非单次尖峰。惊奇触发注意只是工程策略，是否写入长期结构还取决于价值、风险和证据质量。

这些量可能耦合，却要保持解释边界。丰富度增加会提高验证负担，一致性约束过强会压制探索，惊奇阈值过低会造成频繁重构。系统因此求解受限多目标问题，而不是把三者相加成“认知能量”。例如研究助手面对新论文时，可以先以低成本生成宽候选，再用引用核验提高一致性，最后只把经多源支持的高惊奇结果写入结构记忆。若一篇文本词汇宽广但引用错误，其覆盖分高、证据分低；若结果与模型不同但测量仪器故障，其惊奇高、结构更新资格低。分量化指标能够解释这些差异。

验证时，我们构造同义重复、逻辑顺滑但事实错误、正确却表达破碎、低概率传感噪声和真正机制变化等对照样本。指标若把同义重复当丰富、把流畅当正确或把噪声当结构突破，就需要修订。跨模型比较必须使用相同任务、事件划分和标注协议；否则数值大小没有共同尺度。HSF-HD 在此提供的是测量设计原则：状态量必须对应可重复操作，物理词只在确有信号处理对象时使用，任何指标都不能越过其定义域承担本体结论。

在运行时，CCM 可以把覆盖度、冲突率、残差和验证成本作为负荷信号；TCE 根据任务阶段调整探索与核验份额；SPC 选择与当前自我和目标相关的摘要；Offloading 把证据、推导和大规模候选移至外部记忆，减轻 $h(t)$。这套映射保留了“内容、秩序与新异性共同塑造状态”的洞见，同时把每一项变为可以记录、扰动和复现实验的计算量。

\section{语义迁移：跨语境对齐、接口补偿与路径依赖}
\label{sec:geometric_transport}

当系统把“日常生活中的增长”用于数学模型，或把视觉对象交给语言规划器时，移动的不是一个保持绝对形状的语义粒子，而是一组身份、关系和读出条件经过接口转换。设源语境状态空间为 $\mathcal X_A$，目标语境为 $\mathcal X_B$，迁移器为 $T_{A\to B}$。它至少要映射身份、类型、单位、约束和来源：
\begin{equation}
 X_B=T_{A\to B}(X_A;\Gamma_{AB}),
\end{equation}
其中接口契约 $\Gamma_{AB}$ 声明哪些字段可保留、哪些需重编码、哪些无法表达。理解不是无损平行移动，而是在任务允许误差内保持指定查询的答案：对查询集 $Q$，要求 $d_q(D_q^B(T_{A\to B}X),D_q^A(X))\leq\varepsilon_q$。没有查询集和误差尺度，“意义不变”无法被检验。

接口补偿对应原章联络系数想要解决的问题。不同语境可能采用摄氏度与华氏度、对象中心与场景中心、概率与等级标签；迁移器必须应用单位换算、角色重绑定和置信校准。若只复制数值，表示看似相同却语义错误。组合接口还要满足近似一致性：从 $A$ 经 $B$ 到 $C$ 的结果应与直接从 $A$ 到 $C$ 在相关查询上接近。偏差过大说明中间接口丢失了身份、精度或约束，不能用“相位旋转”掩盖。对于不可逆摘要，应明确丢失字段，并保留回到原证据的引用。

路径依赖则是真实而重要的。若系统先学习统计因果再学习社会规范，与相反顺序可能形成不同的默认解释、缓存和技能组合；闭环返回同一主题时，状态不必回到起点。但这一差异来自事件写入、参数更新、资源消耗和承诺变化，不必称为非阿贝尔和乐或贝里相位。我们可定义回路差异
\begin{equation}
 \Delta_\gamma(X)=d\bigl(U_{e_k}\circ\cdots\circ U_{e_1}(X),X\bigr),
\end{equation}
其中距离 $d$ 只比较预先指定的身份、读出或行为。若更新算子不交换，改变事件顺序会改变结果；这是可通过回放两个顺序验证的计算性质。

创造性可以被解释为路径依赖的一种建设性利用：系统跨越多个语境，保留共同身份并组合此前未连接的约束，最后得到新的可检验候选。新颖不等于有效。一个隐喻在文学语境有启发性，映射到工程语境后仍需检查单位、因果和执行条件。以“认知空气动力学”为例，它提示我们研究状态、扰动、控制和稳定边界，但不能直接从流体方程推出智能机制；真正迁移的是问题结构和分析方法，而不是物理量一一同构。好的迁移记录保留对应项与断裂点，使读者知道类比在哪一步结束。

跨语境迁移还需要处理社会来源。不同主体对同一词语可能有不同定义、权限和承诺，系统不能把向量对齐当作共识。接口应保存说话者、时间、制度范围和争议状态；若目标语境没有相应授权，内容可以被理解但不能被执行。外部工具返回的数据也要经过模式验证和版本检查。AgentOS 的 Boundary 可控制哪些字段跨边界，$M_e$ 保存原材料，$E$ 记录迁移事件，$\Theta$ 保存稳定转换规则；一般系统只需实现等价的契约与审计能力。

验证迁移器时，我们使用往返测试、组合路径测试、反事实替换和任务性能测试。往返接近不保证事实正确，但能发现编码丢失；直接与间接路径接近不保证全局一致，但能定位接口偏差；替换身份或单位可以检查系统是否真正使用了契约；最终任务测试确认迁移后的表示是否支持目标操作。由此，“理解”成为在有限查询下保持身份、约束和可用性的能力，而不是不可观测的频率校准。

\section{不变量与学习：身份连续、资源收支和受控重构}
\label{sec:noether_semantion_dynamics}

智能状态在变化中确实需要保持某些量，但不变量来自系统设计、任务约束和更新规则，而不是因为写下一个 $U(1)$ 对称拉格朗日量就自动得到“语义荷守恒”。我们把不变量分为三类。第一类是身份连续性：对象经过属性变化仍能由版本链追踪；第二类是结构约束：类型、权限或事务一致性在允许更新下保持；第三类是资源收支：容量、调用次数和外部预算按真实单位守恒或可核算。信息可以复制、压缩和丢失，语义也会因新证据修订，因此不存在普遍不生不灭的语义总量。

设状态更新为 $X_{t+1}=F_t(X_t,u_t)$，候选不变量为 $I:\mathcal X\to\mathbb R^k$。只有在声明的输入集合 $\mathcal U_0$ 和有效域 $\mathcal D$ 内满足
\begin{equation}
 I(F_t(X,u))=I(X),\qquad X\in\mathcal D, u\in\mathcal U_0,
\end{equation}
才能称 $I$ 在该更新下保持。若外部输入创建新对象、授权者撤销承诺或系统执行遗忘，等式通常不成立，应使用带收支项的平衡式 $I(X_{t+1})-I(X_t)=S_t-L_t$。这类似连续性方程的形式，但 $S_t$ 和 $L_t$ 是可记录的创建与删除事件，不是认知场的物理源汇。逐事件核账比宣称全局守恒更适合审计。

原章的“语义子波包”可以重建为带身份的消息或状态包沿网络传播。设消息 $m=(i,p,c,s,v)$ 包含身份、载荷、上下文、来源和版本。边接口 $T_e$ 作用后得到 $m'=T_e(m)$；若接口承诺保持身份，则必须验证 $i'=i$，同时允许载荷因摘要或单位转换而变化。路径 $\gamma=e_1\cdots e_k$ 上的传播是算子组合 $T_\gamma$。传播成功并不意味着内容零损失，系统应累计误差界、丢失字段与验证结果。若某一步无法满足契约，消息进入隔离或请求补充，而不是依靠有效质量和洛伦兹力继续滑行。

学习是状态传播对结构和参数的受控反作用。令快速状态为 $z_t$，慢参数为 $\theta_t$，验证损失为 $L(\theta;D_t)$。一种离散更新为
\begin{equation}
 \theta_{t+1}=\Pi_{\Theta}\!\left(\theta_t-\eta_t\widehat\nabla L(\theta_t;D_t)\right),
\end{equation}
其中 $\eta_t$ 是步长，$\Pi_{\Theta}$ 把结果投影到允许参数域，梯度估计的偏差和方差需要单独控制。若结构图也改变，必须在编辑后重新定义损失与读出，不能把固定维度的梯度直接延续。一步损失下降不保证长期泛化，更不保证外部任务成功；保留集、分布外样本和真实执行反馈分别检验这些性质。

快慢耦合还需要防止自我强化。错误候选若因高活动被写入结构，下一轮结构又提高该候选活动，就会形成回声回路。控制方法包括来源去重、延迟提交、反证检索、更新上限和独立验证集。对重要结构，系统保存旧版本与回滚条件；对外部已经执行的动作，内部回滚不能撤销现实后果，只能产生补救事件。耗散结构理论在这里提供开放系统需要持续输入、反馈和排出的原则，但具体稳定性仍由更新映射、延迟、增益和资源边界决定。

HSF-HD 因而把整体智能描述为状态、接口与更新规则共同形成的动力系统：局部结构网络维持身份与关系，全局分布表示提供候选与不确定性，规范状态给出目标与边界，读出接口把有限内容提交给行动，反馈再修订快状态和慢结构。AgentOS 中，$h(t)$、$E$、$\Theta$ 分别承担快速切片、事件区间和长期生成结构，$\Psi$、SPC、TCE、CCM、Boundary、Offloading 与 $M_e$ 协调自我相关读出、增益、负荷、边界和外部支撑。它们展示一种实现路径，却不把工程组件提升为所有智能系统的先验公理。

本章的结论不是“思考是弯曲空间中的物理波”，而是：可工作的智能必须把异构状态联合起来，在变化中保持必要身份，在语境转换中声明损失，在有限接口上审查并提交候选，再以证据和反馈修订自身。相位、波包、平行移动和守恒流可以继续作为局部数学工具或辅助意象，但只有对象、条件、单位和验证方法完整时，才进入理论证明。
